\documentclass[../../main.tex]{subfiles}
\begin{document}

\begin{center}
  {\fontsize{48pt}{58pt}\selectfont \textbf{1984}\normalsize\textbf{年度}}
\end{center}

\begin{table}[htb]
  \centering
  \caption{出題テーマと難易度}
  \begin{tabular}{|c|p{6.2cm}|c|c|c|}
    \hline
    \textbf{問題番号} & \textbf{テーマ} & \textbf{思考量} & \textbf{計算量} & \textbf{総合} \\
    \hline
    第1問 & 双曲線を保つ一次変換の条件 &  &  &  \\
    \hline
    第2問 & 合成三角関数の周期 &  & 普 &  \\
    \hline
    第3問 & 動く直線と線分が通過する範囲 & 普 & 普 & 普 \\
    \hline
    第4問 & 内心の位置ベクトルと係数 & 難 & 普 & 普 \\
    \hline
    第5問 & 赤球を先に引く交互ゲームの勝率 & 難 & 難 & 難 \\
    \hline
    第6問 & 導関数に上界がある写像の不動点 & 易 & 易 & 易 \\
    \hline
  \end{tabular}
\end{table}

\begin{table}[htb]
  \centering
  \caption{解答}
  \renewcommand{\arraystretch}{1.4}
  \begin{tabular}{|c|c|p{9cm}|}
    \hline
    \textbf{問題番号} & \textbf{小問} & \textbf{答え} \\
    \hline
    第1問 & -- &  \\
    \hline
    第2問 & (1) & $\sin(\sin x)$の周期は$2\pi$ \\
    \cline{2-3}
    & (2) & $\cos(\sin x)$の周期は$\pi$ \\
    \cline{2-3}
    & (3) & $\sin(x^3)$は周期関数ではない \\
    \hline
    第3問 & (1) & $y\le\dfrac{x^2}{4}$（境界を含む；解答の図を参照） \\
    \cline{2-3}
    & (2) & $\dfrac{7}{6}$ \\
    \hline
    第4問 & (1) & $u=v=w=0$のみ \\
    \cline{2-3}
    & (2) & $(u,v,w)=\dfrac{1}{a+b+c}(a,b,c)$ \\
    \hline
    第5問 & (1) & $P_i\ge P_{i+1}$ \\
    \cline{2-3}
    & (2) & $P_A\ge\dfrac{1}{2}$．等号は$r=1$かつ$s$が奇数のとき \\
    \hline
    第6問 & (1) & 不動点$c_0$がただ1つ存在する \\
    \cline{2-3}
    & (2) & $|d_n-c_0|\le k^n|d_0-c_0|$ \\
    \hline
  \end{tabular}
\end{table}

\paragraph{試験形態}
6問．試験時間は要確認．

\end{document}
